% !TEX root = ../main.tex
\section{Discussão}
\label{sec:discussao}

\subsection{O achado central: parcimônia supera profundidade}

O resultado mais robusto deste trabalho não é o $R^2$ = 0,53 da CNN 2D --- é a
constatação, convergente por duas vias independentes (ablação de dimensionalidade,
Seção~\ref{sec:ablation}; e importância de variáveis do modelo interpretável,
Seção~\ref{sec:interpret}), de que \textbf{a distribuição de altura do dossel,
resumida em poucas dezenas de números, carrega quase todo o sinal preditivo de AGB
disponível nesta amostra}. Arquiteturas de maior capacidade representacional
(PointNet sobre nuvem crua, CNN 3D sobre voxels) não convertem essa riqueza
geométrica adicional em desempenho preditivo superior --- pelo contrário, a
PointNet é consistentemente a pior de todas, e a CNN 3D só se torna competitiva
depois de um controle de qualidade de dados agressivo. Este é o enquadramento
correto da contribuição do trabalho: não ``qual arquitetura é a melhor para AGB por
LiDAR'', mas \textbf{em que ponto a complexidade do modelo deixa de ajudar e passa
a atrapalhar, dado um orçamento de poucas centenas de parcelas rotuladas} --- um
regime realista para a maior parte dos estudos de campo na Amazônia, onde
inventários de centenas de parcelas já representam um esforço logístico
considerável.

\subsection{A limitação mais importante: validação espacialmente ingênua}
\label{sec:spatial-validation}

A validação cruzada k-fold aleatória usada em todos os resultados da
Seção~\ref{sec:resultados} particiona \textbf{parcelas}, não \textbf{sítios}. É
importante separar dois mecanismos distintos pelos quais isso pode inflar o $R^2$
reportado --- eles têm força de evidência diferente, e merecem tratamento
diferente.

\textbf{(a) Pseudo-replicação dentro de uma mesma quadra/campo --- mecanismo que
não depende de heterogeneidade biológica entre sítios.} Muitas das 325 parcelas do
conjunto final não são amostras independentes de ``a floresta amazônica'': são
\textbf{subdivisões contíguas de um mesmo campo original}. Dos 325, \textbf{125
(38\%) pertencem a um grupo de $\geq$2 quadras-irmãs} derivadas do mesmo
\texttt{Name} de campo (ex.: CAU tem campos como T17--T22 cada um dividido em 4--5
quadras adjacentes; FST tem 4 campos divididos em 5 quadras cada) --- e \textbf{um
único sítio (CAU) responde por 85 das 325 parcelas (26\% da amostra inteira)}.
Quadras vizinhas de um mesmo campo compartilham, literalmente, o mesmo dossel
contínuo ou dosséis a poucas dezenas de metros um do outro: seus vetores de altura
são amostras altamente correlacionadas da mesma realização física de floresta, não
realizações independentes de ``estrutura típica da Amazônia''. Um modelo pode
ajustar o padrão numérico específico de um campo (treino) e prever bem sua
quadra-irmã (teste) sem ter aprendido nada que se transporte para um hectare
genuinamente não amostrado --- isso é verdade \textbf{mesmo sob a hipótese de que
toda a floresta amazônica fosse biologicamente homogênea}, porque o problema não é
``os sítios são diferentes entre si'', é ``as unidades de teste e treino, dentro de
um sítio, não são estatisticamente independentes''. Esse mecanismo é o mais sólido
dos dois e o que mais pesa nesta ressalva.

\textbf{(b) Confundidor de campanha/densidade de varredura --- mecanismo
condicional, cuja força depende de quanto os sítios de fato diferem entre si.} Se a
floresta amostrada for de fato aproximadamente homogênea entre sítios (mesmo bioma,
mesma faixa de biomassa típica, sem efeito sistemático de solo/perturbação/região
correlacionado com o alvo) e nenhuma feature identificar o sítio diretamente, este
segundo mecanismo é fraco --- e a crítica de que ``misturar sítios é um problema''
perde força nesse cenário, como argumentado. Mas a premissa de homogeneidade não é
demonstrada pelos próprios dados do projeto: o pipeline de alometria
(\texttt{config/allometry.yaml}) usa curvas altura-diâmetro \textbf{regionais
distintas} (Amazônia centro-oriental, ocidental, Escudo das Guianas, Escudo
Brasileiro) precisamente porque a relação estrutura$\leftrightarrow$biomassa varia
por região --- e a feature $\log(\text{n}^{\text{o}}\text{ de pontos})$ /
\texttt{pt\_density}, presente em todos os modelos, captura densidade de varredura,
que correlaciona com campanha de voo e, por acaso, com região. Este mecanismo
permanece \textbf{plausível, não testado} (Seção~\ref{sec:confounders}) --- mais
fraco que (a), mas não eliminado.

A consequência prática: mesmo que (b) se revele desprezível após a análise
estratificada por campanha (Seção~\ref{sec:confounders}, ainda pendente), o
mecanismo (a) por si só já é suficiente para que os números da
Seção~\ref{sec:resultados} devam ser lidos como ``quanto da variância é explicada
quando se permite reaproveitar informação de quadras vizinhas do mesmo campo'' ---
não ``quão bem isto generaliza para um hectare novo da Amazônia'', que é a pergunta
que a motivação do trabalho (Seção~\ref{sec:intro}) coloca. Fechar essa lacuna
exige reportar uma validação agrupada por sítio (ou, de forma ainda mais
conservadora, por campo/\texttt{Name} de origem da quadra) como métrica principal,
não como nota de rodapé --- com a ressalva prática de que os 18 sítios do conjunto
final têm tamanhos muito desiguais (de 4 a 85 parcelas), então um
\textit{leave-one-site-out} ingênuo produzirá folds de variância muito heterogênea
e deve ser reportado por sítio, não só agregado. É esperável que o $R^2$ caia sob
esse protocolo mais exigente. Isto não invalida o achado central da
Seção~\ref{sec:discussao} (a relação entre dimensionalidade e desempenho deve se
manter qualitativamente, já que ambos os mecanismos acima afetam todos os modelos
comparados de forma semelhante), mas invalida qualquer leitura dos números
absolutos de $R^2$/RMSE como medida de prontidão operacional.

\subsection{Quantificação de incerteza: intervalo conforme por parcela}
\label{sec:conformal}

A motivação declarada deste trabalho --- contabilidade de carbono, REDD+, mercados
de crédito --- é um domínio em que uma estimativa pontual de AGB sem intervalo de
predição tem utilidade limitada: qualquer aplicação de compensação ou relatório de
emissões exige saber não só o valor central, mas o \textbf{intervalo de confiança}
da estimativa por parcela ou por pixel. As métricas reportadas neste artigo
($R^2$, RMSE, rRMSE) descrevem o comportamento agregado do modelo sobre o conjunto
de teste; nenhuma delas produz, por si só, um intervalo de predição individual.

Um intervalo conforme por split (\textit{split conformal prediction}) foi
calculado sobre os resíduos out-of-fold já salvos de cada modelo (variante sem
outliers, o mesmo conjunto reportado na Seção~\ref{sec:arch}): os resíduos
absolutos de metade das parcelas (calibração, $n \approx$ 162) definem o raio do
intervalo a 90\% de confiança; a outra metade, nunca usada na calibração, mede a
cobertura empírica real. O resultado: raios de $\pm$109~Mg/ha (CNN 2D, o modelo
mais acurado) a $\pm$144~Mg/ha (PointNet, o menos acurado), com cobertura empírica
entre 86\% e 92\% --- próxima da nominal de 90\% em todos os seis modelos, o que
indica que o intervalo conforme é razoavelmente bem calibrado nesta amostra,
apesar do tamanho reduzido da calibração. Um raio de $\pm$109 a $\pm$144~Mg/ha
sobre uma AGB média de $\approx$ 170~Mg/ha (variante sem outliers) é, na prática,
um intervalo largo --- a maior parte da AGB plausível de uma parcela amazônica cai
dentro dele --- mas é a primeira quantificação honesta de quão incerta é cada
estimativa individual, ao contrário de reportar só $R^2$/RMSE agregados. Este
intervalo cobre apenas a incerteza do modelo estrutura$\rightarrow$AGB; não propaga
o erro alométrico da Seção~\ref{sec:biomassa} nem a incerteza de medição de campo
até o alvo de treino --- ambos permanecem como trabalho futuro.

\subsection{O alvo de treino é, ele mesmo, um modelo}
\label{sec:allom-sens}

A AGB de campo usada como alvo de treino não é uma medição direta: é a saída de
uma equação alométrica (Seção~\ref{sec:biomassa}) aplicada a medições de campo
(DBH, altura, espécie), e essa equação tem incerteza própria, documentada na
literatura que a origina (Chave et al., 2014). Este trabalho reporta os resultados
principais (Seção~\ref{sec:resultados}) usando \textbf{M1} (densidade de madeira
uniforme, $\rho$ = 0,6 g/cm\textsuperscript{3}) como alvo --- a escolha mais
simples entre as três fórmulas calculadas na Seção~\ref{sec:biomassa} (M2 por tipo
de árvore, M3 por densidade específica de espécie).

Uma análise de sensibilidade foi conduzida: os sete modelos foram retreinados
contra \textbf{M3} (densidade de madeira específica por espécie) sob o mesmo
protocolo de validação cruzada.

\textbf{Ressalva de leitura}: a coluna $R^2$ (M3) abaixo foi treinada quando a
variante de melhor desempenho conhecida para cada modelo ainda era o pipeline
combinado (``solo+filtro+outliers'', sem gap) --- a Seção~\ref{sec:preproc}
identificou depois, via ablação isolada, que a remoção de outliers SOZINHA (sem
solo/filtro) é a variante líder para cinco dos seis modelos, superando inclusive o
combinado em alguns casos. A comparação M1 vs. M3 abaixo ainda não foi refeita
sobre essas variantes isoladas --- a coluna $R^2$ (M1) foi atualizada para
refletir os novos melhores valores (Seção~\ref{sec:arch}), mas a coluna $R^2$ (M3)
permanece com o valor do pipeline combinado antigo; $\Delta R^2$ é, portanto, um
limite inferior pessimista da perda real sob M3 nos modelos onde a variante
isolada lidera.

\begin{table}[H]
  \centering
  \caption{Sensibilidade do desempenho à fórmula alométrica de referência
  (M1 vs. M3).}
  \label{tab:m1m3}
  \begin{tabular}{lccc}
    \toprule
    \textbf{Modelo} & \textbf{$R^2$ (M1)} &
      \textbf{$R^2$ (M3, pipeline combinado)} & \textbf{$\Delta R^2$} \\
    \midrule
    CNN 2D                 & 0,547 & 0,286 & \textbf{$-$0,261} \\
    Random Forest          & 0,532 & 0,382 & $-$0,150 \\
    GBR alturas completas  & 0,526 & 0,309 & $-$0,217 \\
    ABA                    & 0,514 & 0,360 & $-$0,154 \\
    GBR 32 quantis         & 0,510 & 0,329 & $-$0,181 \\
    CNN 3D                 & 0,455 & 0,327 & $-$0,128 \\
    PointNet               & 0,212 & 0,128 & $-$0,084 \\
    \bottomrule
  \end{tabular}
  \fonte{elaborado pelo autor (2026), a partir das validações cruzadas do
  pipeline de treino sob M1 e M3}
\end{table}

Dois achados, nenhum deles trivial. Primeiro, \textbf{todos os sete modelos perdem
poder preditivo substancial sob M3} ($\Delta R^2$ entre $-$0,08 e $-$0,26) --- a
densidade de madeira específica por espécie não é visível na estrutura do dossel
medida pelo LiDAR (nenhuma das representações usadas aqui recebe identidade
taxonômica como insumo), então parte da variância que M3 introduz no alvo é, para
qualquer um destes modelos, ruído irredutível. Segundo, e mais importante para a
validade do resultado central da Seção~\ref{sec:arch}: \textbf{a escolha da fórmula
muda a ordem dos modelos}. Sob M1, a CNN 2D lidera; sob M3 (na comparação
disponível, pipeline combinado), ela cai para o penúltimo lugar (só acima da
PointNet), e o Random Forest assume a liderança --- a maior queda absoluta de toda
a comparação ($\Delta R^2 \approx -0,26$) é justamente do modelo antes apontado
como o melhor. Isso não invalida o argumento de que arquiteturas parcimoniosas
(ABA, Random Forest, GBR) superam a PointNet --- essa conclusão sobrevive intacta
sob as duas fórmulas ---, mas enfraquece a afirmação mais específica de que a
CNN 2D é a melhor arquitetura desta comparação: esse resultado é sensível à
escolha, metodologicamente defensável mas arbitrária entre três candidatas
igualmente válidas, de qual fórmula alométrica usar como alvo. \textbf{Trabalho
pendente}: retreinar a comparação M3 sobre a variante isolada de outliers (líder
sob M1 para cinco dos seis modelos) para confirmar se a inversão de ordem persiste
nessa variante também.

\subsection{Confundidores de campanha e densidade de varredura}
\label{sec:confounders}

As nuvens de pontos usadas neste trabalho vêm de campanhas de voo distintas ---
anos diferentes, sensores possivelmente distintos, densidades de varredura
variáveis --- e uma das features usadas por praticamente todas as abordagens
($\log(\text{n}^{\text{o}}\text{ de pontos})$ ou \texttt{pt\_density}) captura
diretamente essa densidade de aquisição, não uma propriedade da floresta. Isso
abre a possibilidade de que parte do $R^2$ reportado reflita o modelo aprendendo a
\textbf{identificar a campanha} (que pode estar correlacionada com a região e, por
acaso, com a biomassa típica daquela região) em vez de aprender a relação
estrutura $\leftrightarrow$ biomassa que se pretende generalizar.

Duas análises, conduzidas sobre o ABA na sua variante de melhor desempenho
(outliers isolado --- Seção~\ref{sec:preproc}/\ref{sec:arch}), testam essa
hipótese. \textbf{$R^2$ estratificado por ano de campanha} varia bastante --- de
0,12 (2014, $n$=107) a 0,39 (2013, $n$=120), passando por um valor negativo em
2016 ($n$=9, amostra pequena demais para ser conclusiva --- ver ressalva da
Seção~\ref{sec:conformal} sobre fold sizes desbalanceados). Essa variação é
consistente com a hipótese do confundidor, mas não a confirma isoladamente:
campanhas diferentes também cobrem florestas estruturalmente diferentes, então
parte da variação pode ser heterogeneidade biológica real, não viés de densidade
de varredura. Um \textbf{modelo de controle sem as duas features que capturam
densidade de varredura diretamente} (\texttt{pt\_density}, \texttt{ptdens\_cv})
perde pouco desempenho ($R^2$ 0,514 $\rightarrow$ 0,496, $\Delta R^2 \approx$ 0,02)
--- a maior parte do sinal do modelo completo sobrevive sem essas features, o que
pesa contra (mas não elimina) a hipótese de que o modelo está majoritariamente
``reconhecendo a campanha'' pela densidade de pontos. As demais features
(percentis de altura, deciles de densidade por faixa de altura \texttt{D0}--\texttt{D9},
métricas convolucionais do CHM) também podem carregar correlação residual com a
campanha de forma mais sutil, que esta análise não isola.

\subsection{Qualidade do rótulo e heterogeneidade de área}

A auditoria de dados (Apêndice --- Qualidade de Dados) documenta parcelas de área
muito pequena ($\sim$0,02~ha) gerando AGB por hectare implausivelmente alta ---
acima de 1500~Mg/ha, mais que o dobro do teto fisicamente plausível para floresta
tropical densa ---, um artefato estatístico de dividir uma contagem pequena de
árvores por uma área minúscula. O critério de remoção de outliers usado nesta
versão (área $>$ 0,5~ha OU compacidade $<$ 0,6) mitiga a maior parte do problema: a
área das 305 parcelas remanescentes vai de 0,077~ha a 0,305~ha, não mais os
0,02~ha a 1~ha do conjunto completo. Mesmo nessa faixa mais estreita, uma AGB por
hectare estimada sobre poucas árvores tem variância amostral estruturalmente maior
que a mesma métrica sobre uma parcela de 1~ha --- o que significa que as métricas
agregadas de RMSE/$R^2$ da Seção~\ref{sec:resultados} não são igualmente
confiáveis para todas as parcelas do conjunto, mas isto não é reportado de forma
estratificada.

\subsection{Números da literatura: verificação concluída (três erros corrigidos)}
\label{sec:lit-check}

A comparação da Seção~\ref{sec:gedi} usava, em versões anteriores deste texto,
valores extraídos de fontes primárias por pesquisa automatizada sem verificação
adversarial linha a linha. Essa verificação foi concluída e encontrou três erros,
já corrigidos no texto atual: (i) Bispo et al. (2020) estava atribuído à Amazônia
com $R^2$ = 0,58 / rRMSE = 36,0\% --- o estudo é sobre o Cerrado brasileiro e os
valores reais são $R^2$ = 0,93 / rRMSE = 13\% (modelo só-LiDAR); (ii) o DOI citado
para Asner \& Mascaro (2014) apontava para um artigo diferente (albedo boreal,
mesma revista/volume); (iii) o número do GEDI L4A em savana africana
($R^2 \approx$ 0,42, correto) estava atribuído a ``Naidoo et al.'' --- o autor
líder correto é Li et al. (2024), com Naidoo como coautor. Longo et al. (2016) e
Oehmcke et al. (2024) foram confirmados sem ressalva. Esta correção ilustra
exatamente o risco que a versão anterior deste parágrafo já havia antecipado: um
número mal transcrito --- ou, neste caso, três --- compromete a credibilidade de
toda a comparação até ser verificado.

\subsection{Síntese e trabalho futuro}

Em ordem de prioridade para que este trabalho evolua de exploração metodológica
para resultado defensável de dissertação:

\begin{enumerate}
  \item \textbf{Validação LOSO} (Seção~\ref{sec:spatial-validation}) como métrica
    principal, k-fold aleatório rebaixado a referência otimista secundária.
  \item \textbf{Intervalos de predição} (Seção~\ref{sec:conformal}) --- bootstrap
    ou conformal, com cobertura empírica reportada.
  \item \textbf{Análise de sensibilidade alométrica} (Seção~\ref{sec:allom-sens})
    --- resultados da Seção~\ref{sec:resultados} repetidos sob M2 e M3.
  \item \sout{\textbf{Verificação adversarial da literatura}
    (Seção~\ref{sec:lit-check}) antes de qualquer citação em versão final.}
    Concluída --- três erros corrigidos (ver Seção~\ref{sec:lit-check}).
  \item \textbf{Estratificação por campanha e por tamanho de parcela}
    (Seções~\ref{sec:confounders} e seguinte) para descartar confundidores e
    reportar honestamente onde o modelo é, e não é, confiável.
  \item Avaliar formalmente o \textit{ensemble} (média simples de raster + RF +
    voxel, que a exploração preliminar do projeto situa acima de qualquer modelo
    isolado) sob o mesmo protocolo LOSO --- um ganho medido sob k-fold aleatório
    pode não sobreviver ao protocolo mais exigente.
\end{enumerate}

Nenhum desses pontos invalida o achado central da Seção~\ref{sec:discussao}. Pelo
contrário: se a relação entre parcimônia e desempenho se mantiver sob validação
LOSO --- o teste mais rigoroso disponível com os dados atuais ---, ela se torna uma
contribuição metodológica transferível para outros contextos de estimativa de
biomassa por LiDAR com amostras de campo reduzidas, para além do caso amazônico
específico deste trabalho.
